
\UserCommand{Declare Block Type}<\cs{NewBlockType} , \cs{NewBlockType*}>{
    \cs{NewBlockType}
    \{\meta{regular block type}\}
    \AnotherUsage
    \cs{NewBlockType*}
    \{\meta{anonymous block type}\}
}[
    - \meta{regular block type} , \meta{anonymous block type} : The name of the block type to be declared, which can only contain letters and is case-sensitive.
]

You cannot use \cs{NewBlockType} or \cs{NewBlockType*} commands on a block type that has already been created, but you can \seclink{Preamble/Style Configuration}(redefine their styles). \SimpleSystemTeX provides the following predefined basic block types:

\List{Predefined Regular Block Types}{
    - \texttt{Definition} : Mathematical definition.
    - \texttt{Theorem} : Mathematical theorem.
    - \texttt{Corollary} : Mathematical corollary.
}

\List{Predefined Anonymous Block Types}{
    - \texttt{Proof} : Mathematical proof.
}

\newpage
